\documentclass[11pt]{article}
\usepackage[margin=1in]{geometry}
\usepackage{amsmath,amssymb,amsthm}
\usepackage[T1]{fontenc}
\usepackage{lmodern}
\usepackage[hidelinks]{hyperref}

\title{Literature-implied answer: Banach-space subcase of Markov-to-fork convexity}
\author{Run \texttt{fa\_banach\_001}, agent \texttt{agent\_lane\_11}}
\date{2026-06-15}

\newcommand{\dist}{d}
\emergencystretch=2em

\begin{document}
\maketitle

\paragraph{Status.}
\textbf{Literature-implied answer (Banach-space subcase).}
This note records a scoped consequence of known literature for Problem
\texttt{pb:markov->fork} in Baudier--Gartland, \emph{Umbel convexity and the
geometry of trees}, arXiv:2103.16011. It does not solve the problem for
arbitrary metric spaces.

\paragraph{Source question.}
In Section 6, Problem \texttt{pb:markov->fork}, Baudier--Gartland ask whether
Markov \(p\)-convexity implies fork \(p\)-convexity. Immediately before the
problem they explain that Markov \(p\)-convexity averages over all tree levels,
whereas fork \(p\)-convexity involves dyadic-level terms.

\paragraph{Supporting theorem.}
Mendel--Naor, \emph{Markov convexity and local rigidity of distorted metrics},
arXiv:0803.1697, Theorem 1.3, prove that a Banach space is Markov
\(p\)-convex if and only if it admits an equivalent norm whose modulus of
uniform convexity has power type \(p\). Baudier--Gartland cite this result as
Theorem \texttt{thm:MN-LNP}.

\paragraph{Identification.}
Let \(X\) be a Banach space and let \(p\ge 2\). If \(X\), with its original
norm metric, is Markov \(p\)-convex, then by Mendel--Naor it has an equivalent
\(p\)-uniformly convex norm. Baudier--Gartland recall the \(p\)-fork point
inequality in their equation \texttt{eq:qfork}, and in Section 5 note, citing
Mendel--Naor Lemma 2.3, that a \(p\)-uniformly convex Banach norm satisfies the
\(p\)-fork inequality. The full \(p\)-fork point inequality implies the relaxed
fork inequality used in their Theorem \texttt{thm:fork}, hence the equivalent
norm metric is fork \(p\)-convex. Fork \(p\)-convexity is a bi-Lipschitz
invariant, so passing back to the original equivalent Banach norm preserves
the conclusion.

\paragraph{Conclusion.}
For Banach spaces in the renorming range \(p\ge 2\), Markov \(p\)-convexity
does imply fork \(p\)-convexity. The general metric-space question in
Baudier--Gartland remains open and unclaimed here. The supporting theorem
predates the source problem; this packet records an agent-identified subcase,
not an explicit later answer by the supporting authors.

\paragraph{References.}
\begin{itemize}
\item F. P. Baudier and C. Gartland, \emph{Umbel convexity and the geometry of
trees}, arXiv:2103.16011.
\item M. Mendel and A. Naor, \emph{Markov convexity and local rigidity of
distorted metrics}, arXiv:0803.1697; J. Eur. Math. Soc. 15 (2013), 287--337.
\end{itemize}

\end{document}
